\documentclass[10pt,a4paper]{amsart}
\usepackage[margin=19mm]{geometry}
\usepackage{amsmath,amssymb,mathtools}
\usepackage[scr=rsfs,cal=euler]{mathalfa}
\usepackage{xfrac}
\usepackage[unicode,hidelinks,bookmarks=false]{hyperref}
\newcommand{\ie}{i.e.\ }
\newcommand{\cf}{cf.\ }
\newcommand{\resp}{respectively\ }
\newcommand{\Subsection}{\subsection}
\providecommand{\nobreakdash}{\nobreak\hbox{-}}
\setlength{\emergencystretch}{2em}
\multlinegap0pt
\usepackage{amssymb}
\usepackage[normalem]{ulem}
\usepackage{tikz}
\usepackage{cancel}
\newtheorem{lemma}{Lemma}
\title{Conformal boundary rigidity from null geodesic travel times}
\author{Gabriel Paternain, Eric Woolgar}
\date{Source: arXiv:2606.28689v3 (CC BY 4.0)}
\begin{document}
\maketitle

\noindent\emph{Source and licence.} Conformal boundary rigidity from null geodesic travel times. Version arXiv:2606.28689v3 (CC BY 4.0), distributed by arXiv under the Creative Commons Attribution 4.0 International licence, http://creativecommons.org/licenses/by/4.0/, as recorded in the arXiv licence metadata for this exact version and shown on the abstract page https://arxiv.org/abs/2606.28689v3. Full licence notice: \texttt{licenses/arxiv-2606.28689-CC-BY-4.0.txt}.\par\smallskip

\noindent\textit{Editorial scope.} Counted unit: the article's lemma lemma3.2 (source0.tex lines 251--253). The article's complete original proof of it is delivered with it (source0.tex lines 960--1002, one \texttt{proof} environment of 43 lines, which the article places away from the statement). No mathematics is written, repaired or re-derived: the statement and the proof are the article's own lines, in the article's own notation, and no retained line is substituted or re-flowed.  No further source-local context is needed: every cross-reference of every retained line resolves inside the two delivered spans.  Nothing was omitted from any retained span, so the package carries no omission record.  No retained line contains a citation, so the wrapper carries no bibliography and no bibliography entry.  No retained line contains a figure environment, an image-inclusion command or drawing code, so no image is delivered and the package declares no asset dependency.  Editorial formatting only, in addition: an AMS \texttt{amsart} preamble that restates the article's own declarations and macros used by the retained lines, namely the environments \texttt{lemma} on separate counters, together with the standard packages those lines need.  The retained environments print in this document as Lemma 1 in order of appearance, whereas the article prints the same items with the numbering of its own full document.  The complete native source of the article as distributed in the arXiv e-print for this exact version is bundled unchanged as \texttt{sources/203871/source0.tex}.
\begin{lemma}[Polarization lemma]\label{lemma3.2}
Let $T$ be a symmetric $(0,2)$-tensor (at a point) such that $T(\ell,\ell)=0$ for every future-null vector (at that point). Then $T$ is proportional to $g$.
\end{lemma}

\begin{proof}[Proof of Lemma \ref{lemma3.2}]
Fix an orthonormal basis $\left \{ e_0,e_1,\dots,e_{n-1}\right \}$ such that $e_0$ is timelike. Then $e_0\pm e_1$ is null, so
\begin{equation}
\label{eqB.1}
\begin{split}
0=&\, T(e_0+e_1,e_0+e_1)= T(e_0,e_0)+T(e_1,e_1)+2T(e_0,e_1),\\
0=&\, T(e_0-e_1,e_0-e_1)= T(e_0,e_0)+T(e_1,e_1)-2T(e_0,e_1).
\end{split}
\end{equation}
Subtracting one of these from the other, we get
\begin{equation}
\label{eqB.2}
T(e_0,e_1)=0,
\end{equation}
and thus
\begin{equation}
\label{eqB.3}
T(e_0,e_0)=-T(e_1,e_1)
\end{equation}
Replacing $e_1$ by $e_2$ and so forth up to $e_{n-1}$, we see that
\begin{equation}
\label{eqB.4}
\begin{split}
T(e_0,e_i)=&\, 0\text{ for }i=1,2,\dots,n-1,\text{ and }\\
T(e_0,e_0)=&\, -T(e_1,e_1)=-T(e_2,e_2)=\dots=-T(e_{n-1},e_{n-1}).
\end{split}
\end{equation}

Next, $e_0+\frac{1}{\sqrt{2}}\left ( e_1+e_2\right )$ is also null, so (using that $T(e_0,e_1)=0$ and $T(e_0,e_2)=0$) we have that
\begin{equation}
\label{eqB.5}
\begin{split}
0=&\, T(e_0,e_0)+\frac12 T(e_1+e_2,e_1+e_2)\\
=&\, T(e_0,e_0)+\frac12 \left [T(e_1,e_1)+T(e_2,e_2) +2T(e_1,e_2) \right ].
\end{split}
\end{equation}
But $T(e_0,e_0)=-T(e_1,e_1)=-T(e_2,e_2)=-\frac12 \left [ T(e_1,e_1)+T(e_2,e_2)\right ]$, so \eqref{eqB.4} yields $T(e_1,e_2)=0$. Cycling through the basis vectors as before, we have
\begin{equation}
\label{eqB.6}
T(e_i,e_j)=0\quad \forall i,j\in \{ 1,\dots,n-1\}\text{ such that }i\neq j.
\end{equation}
Then from \eqref{eqB.4} and \eqref{eqB.6} we have that $T$ is proportional to $g$.
\end{proof}

\end{document}
